\sectionkicker{Theme synthesis}
\subsection*{横断比較 — Agent能力はどの層へ移ったか}
\addcontentsline{toc}{subsection}{横断比較 — Agent能力はどの層へ移ったか}
2025年後半のagent系更新は、単にbase modelの推論能力を上げる話ではなく、browser、code/tool execution、context、SDK、評価・改善ループ、service lifecycleへ実行面を広げた。ここでは各社の更新を同じ製品比較に潰さず、能力が置かれた層を分けて読む。
\medskip
\begin{center}
\small
\renewcommand{\arraystretch}{1.16}
\begin{tabularx}{\linewidth}{@{}p{0.20\linewidth}X@{}}
\toprule
観察軸 & 一次資料・論文から読めること \\
\midrule
user-facing action surface & ChatGPT agentはresearch、browser automation、code/tool executionを一つのuser-facing productへまとめ、同日にsystem cardも公開した。action surfaceとcontrol boundaryが同じrelease packageで前景化した。 \cite{src-7a7213b68631,src-b074ccdb2d9e} \\
modelとagent platformの結合 & AnthropicのSonnet 4.5、Haiku 4.5、Opus 4.5はAPI提供だけでなくClaude CodeやAgent SDKなどの利用面と結び付いて展開された。living newsroomのため、半年の順序はitem-level chronologyとして扱う。 \cite{src-f6aa679babd7} \\
service lifecycle & Gemini APIではvideo/image、robotics、Computer Use、Gemini 3、Interactions APIなどがpreview・GA・更新として並ぶ。これらはupstream model announcementと同一の日付・identityへ統合しない。 \cite{src-a24fd97eb3d0} \\
build-evaluate-improve loop & AgentKit、Evals、RFT for agentsはagent構築だけでなく評価と改善を反復するengineering layerを示し、Codexの更新はterminal・IDE・web・mobileをまたぐcoding-agent surfaceを広げた。 \cite{src-0661418088b8,src-6deeebce1719} \\
\bottomrule
\end{tabularx}
\renewcommand{\arraystretch}{1.0}
\end{center}
\normalsize
\begin{claimboundary}[この比較の境界]
この表は本号で既に検証・参照している一次資料と論文を横断比較の形へ再配置したものである。vendor / project / author が報告した性能値や優位性は独立再現済みの一般則へ変換せず、本文とSource Notesの帰属境界をそのまま維持する。
\end{claimboundary}
